% Part: normal-modal-logic
% Chapter: filtrations
% Section: preliminaries

\documentclass[../../../include/open-logic-section]{subfiles}

\begin{document}

\olfileid{nml}{fil}{pre}

\olsection{पूर्वतयारी}

निस्यंदनांमुळे आपल्या मोडल तर्कशास्त्रांच्या
प्रणाली निर्णेय आहेत हे सिद्ध करता येते:
त्यांना \emph{सांत प्रतिमान गुणधर्म}
आहे हे त्याने दाखवता येते. म्हणजे,
एखाद्या प्रतिमानात सत्य (असत्य)
असलेले कोणतेही !!{formula}
एखाद्या \emph{सांत} प्रतिमानातही
सत्य (असत्य) असते. निस्यंदनांची
व्याख्या उपसूत्रांखाली बंद असलेल्या
!!{formula}s च्या संचांच्या
सापेक्ष केली जाते.

\begin{defn}\ollabel{defn:modallyclosed}
  !!{formula}s चा संच~$\Gamma$
  हा \emph{उपसूत्रांखाली बंद}
  असेल, जर त्यात~$\Gamma$
  मधील प्रत्येक !!a{formula}
  चे प्रत्येक उपसूत्र असेल.
  शिवाय, $\Gamma$ हा
  \emph{मोडलदृष्ट्या बंद}
  असेल, जर तो उपसूत्रांखाली
  बंद असेल आणि
  $!A \in \Gamma$ असल्यास
  $\Box!A,
  \Diamond!A \in \Gamma$
  हेही खरे असेल.
\end{defn}

उदाहरणार्थ, !!a{formula}~$!A$
दिल्यावर त्याच्या सर्व
उप-!!{formula}s चा संच
उप-!!{formula}s खाली बंद
असतो. $!A$ च्या उप-!!{formula}s
च्या संचातून एखाद्या प्रतिमानाचे
निस्यंदन परिभाषित करताना
हवा असलेला गुणधर्म त्याला
असेल: मूळ प्रतिमानात~$!A$
सत्य (असत्य) असेल तेव्हा
आणि तेव्हाच निस्यंदनातही
ते सत्य (असत्य) असेल.

प्रतिमान~$\mModel{M}$ चे
$\Gamma$ मधून मिळणाऱ्या
निस्यंदनातील जगांचा संच
पुढील तुल्यता संबंधाच्या
सर्व तुल्यतावर्गांचा संच
म्हणून परिभाषित होतो.

\begin{defn}
$\mModel{M} =\tuple{W, R, V}$
असू द्या आणि $\Gamma$ हा
उप-!!{formula}s खाली बंद
आहे असे समजा.
$\Gamma$ मधील तीच
!!{formula}s सत्य ठरवणाऱ्या
कोणत्याही दोन जगांमध्ये
संबंध~$\equiv$ लागू होईल
अशी~$W$ वर त्याची व्याख्या
करू; म्हणजे:
\[
u \equiv v \quad \text{तेव्हा आणि तेव्हाच} \quad
\forall !A \in \Gamma : \mSat{M}{!A}[u] \Leftrightarrow \mSat{M}{!A}[v].
\]
जग~$w$ चा तुल्यतावर्ग
$[w]_\equiv$, किंवा संक्षेपाने
$[w]$, म्हणजे~$w$ शी
$\equiv$-समतुल्य असलेल्या
सर्व जगांचा संच:
\[
[w] = \Setabs{v}{v \equiv w}.
\]
\end{defn}

\begin{prop}
  $\mModel{M}$ आणि $\Gamma$
  दिल्यावर वरीलप्रमाणे
  परिभाषित $\equiv$ हा
  तुल्यता संबंध आहे;
  म्हणजे तो परावर्ती,
  सममित आणि संक्रमक आहे.
\end{prop}

\begin{proof}
  $\equiv$ हा परावर्ती आहे,
  कारण~$w$ येथे सत्य असलेली
  $\Gamma$ मधील !!{formula}s
स्वतः येथील सूत्रांशी
नेमकी जुळतात.
  तो सममित आहे: $u$ आणि~$v$
  येथे $\Gamma$ मधील तीच
  !!{formula}s सत्य असतील,
  तर~$v$ आणि~$u$
  यांच्याबाबतीतही तसेच.
  तो संक्रमकही आहे:
  $u$ आणि~$v$ येथे
  $\Gamma$ मधील तीच
  !!{formula}s सत्य असतील,
  आणि~$v$ व~$w$ येथेही
  तीच असतील, तर~$u$
  व~$w$ येथेही
  $\Gamma$ मधील तीच
  !!{formula}s सत्य असतील.
\end{proof}

इतर प्रत्येक तुल्यता
संबंधाप्रमाणे $\equiv$
हा~$W$ चे \emph{विभाजन}
करतो; म्हणजे~$W$ चे
परस्पर असंयुक्त उपसंच
मिळतात, आणि त्यांचा
संयोग संपूर्ण~$W$ असतो.
प्रत्येक $w \in W$
हा त्यांपैकी एका
विभागाचा, म्हणजे~$[w]$
चा, !!a{element} असतो,
कारण $w \equiv w$.
म्हणून $[w]$ हे विभाग
सर्व~$W$ व्यापतात.
ते परस्पर असंयुक्तही
आहेत: $u \in [w]$
आणि $u \in [v]$
असतील, तर $u \equiv w$
आणि $u \equiv v$.
सममिती व संक्रमकतेवरून
$w \equiv v$; म्हणून
$[w] = [v]$.

\end{document}
